\slajdid{09_3-s002}
\begin{frame}[label=09_3-s002]{Dinamičke karakteristike CMOS invertora}
MOS FET tranzistori po prirodi svog funkcionisanja imaju velike interne kapacitivnosti.
\par\medskip
Na primer, da bi se formirao kanal moraju da se dovedu naelektrisanja na gejt tranzistora koja će iz osnove „privući“ naelektrisanja potrebna za formiranje kanala.
\par\medskip
Ova i slične kapacitivnosti koje poseduje tranzistor će dominantno uticati na brzinu rada.
\par\medskip
Te kapacitivnosti opterećuju sam invertor ali isto tako i prethodni invertor koji mu prosleđuje logičke nivoe.
\par\medskip
Parazitne kapacitivnosti vodova i dalje postoje, ali njih možemo smatrati relativno konstantnim.
\par\medskip
Problem sa kapacitivnostima CMOS invertora je što zavise od dimenzija tranzistora. I dok smo u prethodnim slučajevima menjali dimenzije tranzistora da bi dobili odgovarajuće statičke karakteristike, videćemo da će to itekako uticati na dinamički režim
\end{frame}
